% --- Archon physics formalization source begin ---
% archon:physics
% archon:covers PhyXMiniProblems/problem_phyx_mini_0937.lean
% archon:source-report reports/phyx_mini/problem_phyx_mini_0937.source.json

\chapter{Physics problem phyx\_mini\_0937}
\label{ch:PhyXMiniProblems_problem_phyx_mini_0937}

\paragraph{Problem source.}
\#\# Physical scenario

Suppose the moving rod in figure is \textbackslash{}( 0.10 \textbackslash{}\textbackslash{} \textbackslash{}mathrm\{m\} \textbackslash{}) long, the velocity \textbackslash{}( v \textbackslash{}) is \textbackslash{}( 2.5 \textbackslash{}\textbackslash{} \textbackslash{}mathrm\{m/s\} \textbackslash{}), the total resistance of the loop is \textbackslash{}( 0.030 \textbackslash{}\textbackslash{} \textbackslash{}Omega \textbackslash{}), and \textbackslash{}( B \textbackslash{}) is \textbackslash{}( 0.60 \textbackslash{}\textbackslash{} \textbackslash{}mathrm\{T\} \textbackslash{}).

\#\# Question

Find the force acting on the rod.

\#\# Answer choices

- A: A: 0.45N

- B: B: 0.3N

- C: C: 0.5N

- D: D: 0.2N

\#\# Figure caption (auxiliary; use the image as primary evidence)

The image is a diagram related to electromagnetism, depicting a U-shaped conducting track with a sliding bar. Here's a detailed description:

1. **Conducting Track and Bar**:
   - A U-shaped conductor is shown, with parallel horizontal sections labeled \textbackslash{}(a\textbackslash{}) and \textbackslash{}(b\textbackslash{}).
   - A vertically oriented sliding bar is positioned on the track, capable of moving horizontally.

2. **Magnetic Field**:
   - Blue "X" symbols fill the background area within the U-shaped conductor, indicating a magnetic field directed into the page. This is delineated by the vector \textbackslash{}(\textbackslash{}vec\{B\}\textbackslash{}).

3. **Current and Induced EMF**:
   - There are purple arrows labeled \textbackslash{}(I\textbackslash{}) on the horizontal sections of the conductor, showing the direction of current flow.
   - A blue upward arrow labeled \textbackslash{}(\textbackslash{}mathcal\{E\}\textbackslash{}) represents the induced electromotive force (EMF) across the conductor due to the movement of the bar.

4. **Velocity**:
   - A green arrow labeled \textbackslash{}(\textbackslash{}vec\{v\}\textbackslash{}) points to the right, attached to the sliding bar, indicating the direction of motion.

Overall, the diagram is illustrating a physical scenario involving electromagnetic induction due to the movement of a conductor within a magnetic field.

\#\# Dataset metadata

Electromagnetic Induction; Physical Model Grounding Reasoning, Multi-Formula Reasoning

\paragraph{Recorded answer/context.}
B: B: 0.3N

\paragraph{Figure/image path.}
/root/proposal\_for\_physic/science-mango/phyx\_mini\_run/phyx\_data/test\_image/937.png

\paragraph{Formalization target.}
create a compiling Lean file with sorry bodies at `PhyXMiniProblems/problem\_phyx\_mini\_0937.lean`. The Lean declarations must preserve the physical quantities, dimensions or dimensional roles, figure labels, governing-law hypotheses, and final relation expressed by this problem.
Use Mathlib/Physlib names found through LeanExplore where available. If a physics API is missing, introduce faithful local abstractions rather than scalar placeholder aliases.

\begin{theorem}[Physics formalization target]
\label{thm:physics:phyx_mini_0937:target}
\lean{PhyXMiniProblems.ProblemPhyXMini0937.force_on_moving_rod}
\uses{def:physics:phyx-mini-0937:phyxminiproblems-problemphyxmini0937-forceinnewtons, def:physics:phyx-mini-0937:phyxminiproblems-problemphyxmini0937-slidewireforcesetup, def:physics:phyx-mini-0937:phyxminiproblems-problemphyxmini0937-matchesslidewireforcescenario, def:physics:phyx-mini-0937:phyxminiproblems-problemphyxmini0937-matchesprimaryslidewirefigure, def:physics:phyx-mini-0937:phyxminiproblems-problemphyxmini0937-matchesgivenslidewiremeasurements, def:physics:phyx-mini-0937:phyxminiproblems-problemphyxmini0937-hasphysicalslidewireforceparameters, def:physics:phyx-mini-0937:phyxminiproblems-problemphyxmini0937-hasuniformappliedmagneticfield, def:physics:phyx-mini-0937:phyxminiproblems-problemphyxmini0937-satisfiesslidewireforcelaws}
The assigned autoformalize agent should translate this physics problem into Lean declarations in the covered file, with theorem and lemma proof bodies written as `by sorry`.
\end{theorem}
\begin{proof}
This is an autoformalization task, not a proof task. Produce statements that can later be proved by the physics prover without weakening the physical model.
\end{proof}

% --- Archon declaration topology begin ---
\section{Declaration topology}
\begin{definition}[magnetic Flux Density Dimension]
  \label{def:physics:phyx-mini-0937:phyxminiproblems-problemphyxmini0937-magneticfluxdensitydimension}
  \lean{PhyXMiniProblems.ProblemPhyXMini0937.magneticFluxDensityDimension}
  The dimension M T⁻¹ C⁻¹ of magnetic flux density (tesla)
\end{definition}
\begin{proof}
  This is the declared model definition of magnetic Flux Density Dimension.
\end{proof}
\begin{definition}[electromotive Force Dimension]
  \label{def:physics:phyx-mini-0937:phyxminiproblems-problemphyxmini0937-electromotiveforcedimension}
  \lean{PhyXMiniProblems.ProblemPhyXMini0937.electromotiveForceDimension}
  The dimension M L² T⁻² C⁻¹ of electromotive force (volt)
\end{definition}
\begin{proof}
  This is the declared model definition of electromotive Force Dimension.
\end{proof}
\begin{definition}[electric Current Dimension]
  \label{def:physics:phyx-mini-0937:phyxminiproblems-problemphyxmini0937-electriccurrentdimension}
  \lean{PhyXMiniProblems.ProblemPhyXMini0937.electricCurrentDimension}
  The dimension C T⁻¹ of electric current (ampere)
\end{definition}
\begin{proof}
  This is the declared model definition of electric Current Dimension.
\end{proof}
\begin{definition}[electrical Resistance Dimension]
  \label{def:physics:phyx-mini-0937:phyxminiproblems-problemphyxmini0937-electricalresistancedimension}
  \lean{PhyXMiniProblems.ProblemPhyXMini0937.electricalResistanceDimension}
  The dimension M L² T⁻¹ C⁻² of electrical resistance (ohm)
\end{definition}
\begin{proof}
  This is the declared model definition of electrical Resistance Dimension.
\end{proof}
\begin{definition}[force Dimension]
  \label{def:physics:phyx-mini-0937:phyxminiproblems-problemphyxmini0937-forcedimension}
  \lean{PhyXMiniProblems.ProblemPhyXMini0937.forceDimension}
  The dimension M L T⁻² of force (newton)
\end{definition}
\begin{proof}
  This is the declared model definition of force Dimension.
\end{proof}
\begin{definition}[Magnetic Flux Density Magnitude]
  \label{def:physics:phyx-mini-0937:phyxminiproblems-problemphyxmini0937-magneticfluxdensitymagnitude}
  \lean{PhyXMiniProblems.ProblemPhyXMini0937.MagneticFluxDensityMagnitude}
  \uses{def:physics:phyx-mini-0937:phyxminiproblems-problemphyxmini0937-magneticfluxdensitydimension}
  A nonnegative, unit-independent magnetic-flux-density magnitude
\end{definition}
\begin{proof}
  This is the declared model definition of Magnetic Flux Density Magnitude.
\end{proof}
\begin{definition}[Length Magnitude]
  \label{def:physics:phyx-mini-0937:phyxminiproblems-problemphyxmini0937-lengthmagnitude}
  \lean{PhyXMiniProblems.ProblemPhyXMini0937.LengthMagnitude}
  A nonnegative, unit-independent physical length
\end{definition}
\begin{proof}
  This is the declared model definition of Length Magnitude.
\end{proof}
\begin{definition}[Speed Magnitude]
  \label{def:physics:phyx-mini-0937:phyxminiproblems-problemphyxmini0937-speedmagnitude}
  \lean{PhyXMiniProblems.ProblemPhyXMini0937.SpeedMagnitude}
  A nonnegative, unit-independent speed
\end{definition}
\begin{proof}
  This is the declared model definition of Speed Magnitude.
\end{proof}
\begin{definition}[Electromotive Force Magnitude]
  \label{def:physics:phyx-mini-0937:phyxminiproblems-problemphyxmini0937-electromotiveforcemagnitude}
  \lean{PhyXMiniProblems.ProblemPhyXMini0937.ElectromotiveForceMagnitude}
  \uses{def:physics:phyx-mini-0937:phyxminiproblems-problemphyxmini0937-electromotiveforcedimension}
  A nonnegative, unit-independent electromotive-force magnitude
\end{definition}
\begin{proof}
  This is the declared model definition of Electromotive Force Magnitude.
\end{proof}
\begin{definition}[Electric Current Magnitude]
  \label{def:physics:phyx-mini-0937:phyxminiproblems-problemphyxmini0937-electriccurrentmagnitude}
  \lean{PhyXMiniProblems.ProblemPhyXMini0937.ElectricCurrentMagnitude}
  \uses{def:physics:phyx-mini-0937:phyxminiproblems-problemphyxmini0937-electriccurrentdimension}
  A nonnegative, unit-independent electric-current magnitude
\end{definition}
\begin{proof}
  This is the declared model definition of Electric Current Magnitude.
\end{proof}
\begin{definition}[Resistance Magnitude]
  \label{def:physics:phyx-mini-0937:phyxminiproblems-problemphyxmini0937-resistancemagnitude}
  \lean{PhyXMiniProblems.ProblemPhyXMini0937.ResistanceMagnitude}
  \uses{def:physics:phyx-mini-0937:phyxminiproblems-problemphyxmini0937-electricalresistancedimension}
  A nonnegative, unit-independent electrical resistance
\end{definition}
\begin{proof}
  This is the declared model definition of Resistance Magnitude.
\end{proof}
\begin{definition}[Force Magnitude]
  \label{def:physics:phyx-mini-0937:phyxminiproblems-problemphyxmini0937-forcemagnitude}
  \lean{PhyXMiniProblems.ProblemPhyXMini0937.ForceMagnitude}
  \uses{def:physics:phyx-mini-0937:phyxminiproblems-problemphyxmini0937-forcedimension}
  A nonnegative, unit-independent magnetic-force magnitude
\end{definition}
\begin{proof}
  This is the declared model definition of Force Magnitude.
\end{proof}
\begin{definition}[magnetic Flux Density In Teslas]
  \label{def:physics:phyx-mini-0937:phyxminiproblems-problemphyxmini0937-magneticfluxdensityinteslas}
  \lean{PhyXMiniProblems.ProblemPhyXMini0937.magneticFluxDensityInTeslas}
  \uses{def:physics:phyx-mini-0937:phyxminiproblems-problemphyxmini0937-magneticfluxdensitymagnitude}
  Coherent-SI readout of magnetic flux density, in teslas
\end{definition}
\begin{proof}
  This is the declared model definition of magnetic Flux Density In Teslas.
\end{proof}
\begin{definition}[length In Meters]
  \label{def:physics:phyx-mini-0937:phyxminiproblems-problemphyxmini0937-lengthinmeters}
  \lean{PhyXMiniProblems.ProblemPhyXMini0937.lengthInMeters}
  \uses{def:physics:phyx-mini-0937:phyxminiproblems-problemphyxmini0937-lengthmagnitude}
  Coherent-SI readout of length, in metres
\end{definition}
\begin{proof}
  This is the declared model definition of length In Meters.
\end{proof}
\begin{definition}[speed In Meters Per Second]
  \label{def:physics:phyx-mini-0937:phyxminiproblems-problemphyxmini0937-speedinmeterspersecond}
  \lean{PhyXMiniProblems.ProblemPhyXMini0937.speedInMetersPerSecond}
  \uses{def:physics:phyx-mini-0937:phyxminiproblems-problemphyxmini0937-speedmagnitude}
  Coherent-SI readout of speed, in metres per second
\end{definition}
\begin{proof}
  This is the declared model definition of speed In Meters Per Second.
\end{proof}
\begin{definition}[electromotive Force In Volts]
  \label{def:physics:phyx-mini-0937:phyxminiproblems-problemphyxmini0937-electromotiveforceinvolts}
  \lean{PhyXMiniProblems.ProblemPhyXMini0937.electromotiveForceInVolts}
  \uses{def:physics:phyx-mini-0937:phyxminiproblems-problemphyxmini0937-electromotiveforcemagnitude}
  Coherent-SI readout of electromotive force, in volts
\end{definition}
\begin{proof}
  This is the declared model definition of electromotive Force In Volts.
\end{proof}
\begin{definition}[electric Current In Amperes]
  \label{def:physics:phyx-mini-0937:phyxminiproblems-problemphyxmini0937-electriccurrentinamperes}
  \lean{PhyXMiniProblems.ProblemPhyXMini0937.electricCurrentInAmperes}
  \uses{def:physics:phyx-mini-0937:phyxminiproblems-problemphyxmini0937-electriccurrentmagnitude}
  Coherent-SI readout of electric current, in amperes
\end{definition}
\begin{proof}
  This is the declared model definition of electric Current In Amperes.
\end{proof}
\begin{definition}[resistance In Ohms]
  \label{def:physics:phyx-mini-0937:phyxminiproblems-problemphyxmini0937-resistanceinohms}
  \lean{PhyXMiniProblems.ProblemPhyXMini0937.resistanceInOhms}
  \uses{def:physics:phyx-mini-0937:phyxminiproblems-problemphyxmini0937-resistancemagnitude}
  Coherent-SI readout of resistance, in ohms
\end{definition}
\begin{proof}
  This is the declared model definition of resistance In Ohms.
\end{proof}
\begin{definition}[force In Newtons]
  \label{def:physics:phyx-mini-0937:phyxminiproblems-problemphyxmini0937-forceinnewtons}
  \lean{PhyXMiniProblems.ProblemPhyXMini0937.forceInNewtons}
  \uses{def:physics:phyx-mini-0937:phyxminiproblems-problemphyxmini0937-forcemagnitude}
  Coherent-SI readout of force, in newtons
\end{definition}
\begin{proof}
  This is the declared model definition of force In Newtons.
\end{proof}
\begin{definition}[Rail Label]
  \label{def:physics:phyx-mini-0937:phyxminiproblems-problemphyxmini0937-raillabel}
  \lean{PhyXMiniProblems.ProblemPhyXMini0937.RailLabel}
  The literal labels printed beside the upper and lower contact points
\end{definition}
\begin{proof}
  This is the declared model definition of Rail Label.
\end{proof}
\begin{definition}[Conductor Part]
  \label{def:physics:phyx-mini-0937:phyxminiproblems-problemphyxmini0937-conductorpart}
  \lean{PhyXMiniProblems.ProblemPhyXMini0937.ConductorPart}
  The two conducting pieces in the closed slidewire circuit
\end{definition}
\begin{proof}
  This is the declared model definition of Conductor Part.
\end{proof}
\begin{definition}[Spatial Direction]
  \label{def:physics:phyx-mini-0937:phyxminiproblems-problemphyxmini0937-spatialdirection}
  \lean{PhyXMiniProblems.ProblemPhyXMini0937.SpatialDirection}
  Named side-view directions used by the supplied figure
\end{definition}
\begin{proof}
  This is the declared model definition of Spatial Direction.
\end{proof}
\begin{definition}[opposite]
  \label{def:physics:phyx-mini-0937:phyxminiproblems-problemphyxmini0937-spatialdirection-opposite}
  \lean{PhyXMiniProblems.ProblemPhyXMini0937.SpatialDirection.opposite}
  \uses{def:physics:phyx-mini-0937:phyxminiproblems-problemphyxmini0937-spatialdirection}
  Reversal of a displayed spatial direction
\end{definition}
\begin{proof}
  This is the declared model definition of opposite.
\end{proof}
\begin{definition}[Circuit Sense]
  \label{def:physics:phyx-mini-0937:phyxminiproblems-problemphyxmini0937-circuitsense}
  \lean{PhyXMiniProblems.ProblemPhyXMini0937.CircuitSense}
  Orientation of the current around the loop as seen on the page
\end{definition}
\begin{proof}
  This is the declared model definition of Circuit Sense.
\end{proof}
\begin{definition}[Axis Orientation]
  \label{def:physics:phyx-mini-0937:phyxminiproblems-problemphyxmini0937-axisorientation}
  \lean{PhyXMiniProblems.ProblemPhyXMini0937.AxisOrientation}
  Axis orientation of a drawn conductor
\end{definition}
\begin{proof}
  This is the declared model definition of Axis Orientation.
\end{proof}
\begin{definition}[Slidewire Figure]
  \label{def:physics:phyx-mini-0937:phyxminiproblems-problemphyxmini0937-slidewirefigure}
  \lean{PhyXMiniProblems.ProblemPhyXMini0937.SlidewireFigure}
  \uses{def:physics:phyx-mini-0937:phyxminiproblems-problemphyxmini0937-raillabel, def:physics:phyx-mini-0937:phyxminiproblems-problemphyxmini0937-conductorpart, def:physics:phyx-mini-0937:phyxminiproblems-problemphyxmini0937-spatialdirection, def:physics:phyx-mini-0937:phyxminiproblems-problemphyxmini0937-circuitsense, def:physics:phyx-mini-0937:phyxminiproblems-problemphyxmini0937-axisorientation}
  Literal qualitative content of image 937.png. The raster supplies the labels and directions below but does not itself print the numerical data or a force value
\end{definition}
\begin{proof}
  This is the declared model definition of Slidewire Figure.
\end{proof}
\begin{definition}[Slidewire Force Setup]
  \label{def:physics:phyx-mini-0937:phyxminiproblems-problemphyxmini0937-slidewireforcesetup}
  \lean{PhyXMiniProblems.ProblemPhyXMini0937.SlidewireForceSetup}
  \uses{def:physics:phyx-mini-0937:phyxminiproblems-problemphyxmini0937-magneticfluxdensitymagnitude, def:physics:phyx-mini-0937:phyxminiproblems-problemphyxmini0937-lengthmagnitude, def:physics:phyx-mini-0937:phyxminiproblems-problemphyxmini0937-speedmagnitude, def:physics:phyx-mini-0937:phyxminiproblems-problemphyxmini0937-electromotiveforcemagnitude, def:physics:phyx-mini-0937:phyxminiproblems-problemphyxmini0937-electriccurrentmagnitude, def:physics:phyx-mini-0937:phyxminiproblems-problemphyxmini0937-resistancemagnitude, def:physics:phyx-mini-0937:phyxminiproblems-problemphyxmini0937-forcemagnitude, def:physics:phyx-mini-0937:phyxminiproblems-problemphyxmini0937-spatialdirection, def:physics:phyx-mini-0937:phyxminiproblems-problemphyxmini0937-slidewirefigure}
  Independent observables at one instant. In particular, the magnetic force is not defined from the requested numerical answer
\end{definition}
\begin{proof}
  This is the declared model definition of Slidewire Force Setup.
\end{proof}
\begin{definition}[Matches Slidewire Force Scenario]
  \label{def:physics:phyx-mini-0937:phyxminiproblems-problemphyxmini0937-matchesslidewireforcescenario}
  \lean{PhyXMiniProblems.ProblemPhyXMini0937.MatchesSlidewireForceScenario}
  \uses{def:physics:phyx-mini-0937:phyxminiproblems-problemphyxmini0937-slidewireforcesetup}
  The prose and diagram identify a moving rod closing a U-shaped track
\end{definition}
\begin{proof}
  This is the declared model definition of Matches Slidewire Force Scenario.
\end{proof}
\begin{definition}[Matches Primary Slidewire Figure]
  \label{def:physics:phyx-mini-0937:phyxminiproblems-problemphyxmini0937-matchesprimaryslidewirefigure}
  \lean{PhyXMiniProblems.ProblemPhyXMini0937.MatchesPrimarySlidewireFigure}
  \uses{def:physics:phyx-mini-0937:phyxminiproblems-problemphyxmini0937-slidewireforcesetup}
  Primary-raster evidence, including the literal a and b labels. No force magnitude or answer choice is included here
\end{definition}
\begin{proof}
  This is the declared model definition of Matches Primary Slidewire Figure.
\end{proof}
\begin{definition}[Matches Given Slidewire Measurements]
  \label{def:physics:phyx-mini-0937:phyxminiproblems-problemphyxmini0937-matchesgivenslidewiremeasurements}
  \lean{PhyXMiniProblems.ProblemPhyXMini0937.MatchesGivenSlidewireMeasurements}
  \uses{def:physics:phyx-mini-0937:phyxminiproblems-problemphyxmini0937-magneticfluxdensityinteslas, def:physics:phyx-mini-0937:phyxminiproblems-problemphyxmini0937-lengthinmeters, def:physics:phyx-mini-0937:phyxminiproblems-problemphyxmini0937-speedinmeterspersecond, def:physics:phyx-mini-0937:phyxminiproblems-problemphyxmini0937-resistanceinohms, def:physics:phyx-mini-0937:phyxminiproblems-problemphyxmini0937-slidewireforcesetup}
  Exact coherent-SI measurements stated in the problem text
\end{definition}
\begin{proof}
  This is the declared model definition of Matches Given Slidewire Measurements.
\end{proof}
\begin{definition}[Has Physical Slidewire Force Parameters]
  \label{def:physics:phyx-mini-0937:phyxminiproblems-problemphyxmini0937-hasphysicalslidewireforceparameters}
  \lean{PhyXMiniProblems.ProblemPhyXMini0937.HasPhysicalSlidewireForceParameters}
  \uses{def:physics:phyx-mini-0937:phyxminiproblems-problemphyxmini0937-magneticfluxdensityinteslas, def:physics:phyx-mini-0937:phyxminiproblems-problemphyxmini0937-lengthinmeters, def:physics:phyx-mini-0937:phyxminiproblems-problemphyxmini0937-speedinmeterspersecond, def:physics:phyx-mini-0937:phyxminiproblems-problemphyxmini0937-resistanceinohms, def:physics:phyx-mini-0937:phyxminiproblems-problemphyxmini0937-slidewireforcesetup}
  Positivity and non-vacuity conditions for the modeled physical state
\end{definition}
\begin{proof}
  This is the declared model definition of Has Physical Slidewire Force Parameters.
\end{proof}
\begin{definition}[Has Uniform Applied Magnetic Field]
  \label{def:physics:phyx-mini-0937:phyxminiproblems-problemphyxmini0937-hasuniformappliedmagneticfield}
  \lean{PhyXMiniProblems.ProblemPhyXMini0937.HasUniformAppliedMagneticField}
  \uses{def:physics:phyx-mini-0937:phyxminiproblems-problemphyxmini0937-magneticfluxdensityinteslas, def:physics:phyx-mini-0937:phyxminiproblems-problemphyxmini0937-slidewireforcesetup}
  The scalar flux-density magnitude calibrates the norm of Physlib's magnetic vector field throughout the region marked by uniform crosses
\end{definition}
\begin{proof}
  This is the declared model definition of Has Uniform Applied Magnetic Field.
\end{proof}
\begin{definition}[Satisfies Slidewire Force Laws]
  \label{def:physics:phyx-mini-0937:phyxminiproblems-problemphyxmini0937-satisfiesslidewireforcelaws}
  \lean{PhyXMiniProblems.ProblemPhyXMini0937.SatisfiesSlidewireForceLaws}
  \uses{def:physics:phyx-mini-0937:phyxminiproblems-problemphyxmini0937-magneticfluxdensityinteslas, def:physics:phyx-mini-0937:phyxminiproblems-problemphyxmini0937-lengthinmeters, def:physics:phyx-mini-0937:phyxminiproblems-problemphyxmini0937-speedinmeterspersecond, def:physics:phyx-mini-0937:phyxminiproblems-problemphyxmini0937-electromotiveforceinvolts, def:physics:phyx-mini-0937:phyxminiproblems-problemphyxmini0937-electriccurrentinamperes, def:physics:phyx-mini-0937:phyxminiproblems-problemphyxmini0937-resistanceinohms, def:physics:phyx-mini-0937:phyxminiproblems-problemphyxmini0937-forceinnewtons, def:physics:phyx-mini-0937:phyxminiproblems-problemphyxmini0937-slidewireforcesetup}
  School-physics laws for a rod perpendicular to its velocity and the uniform field. They relate independent observables; none gives the requested 0.3 N value
\end{definition}
\begin{proof}
  This is the declared model definition of Satisfies Slidewire Force Laws.
\end{proof}
\begin{lemma}[magnetic force magnitude formula]
  \label{lem:physics:phyx-mini-0937:phyxminiproblems-problemphyxmini0937-magnetic-force-magnitude-formula}
  \lean{PhyXMiniProblems.ProblemPhyXMini0937.magnetic_force_magnitude_formula}
  \uses{def:physics:phyx-mini-0937:phyxminiproblems-problemphyxmini0937-magneticfluxdensityinteslas, def:physics:phyx-mini-0937:phyxminiproblems-problemphyxmini0937-lengthinmeters, def:physics:phyx-mini-0937:phyxminiproblems-problemphyxmini0937-speedinmeterspersecond, def:physics:phyx-mini-0937:phyxminiproblems-problemphyxmini0937-resistanceinohms, def:physics:phyx-mini-0937:phyxminiproblems-problemphyxmini0937-forceinnewtons, def:physics:phyx-mini-0937:phyxminiproblems-problemphyxmini0937-slidewireforcesetup, def:physics:phyx-mini-0937:phyxminiproblems-problemphyxmini0937-hasphysicalslidewireforceparameters, def:physics:phyx-mini-0937:phyxminiproblems-problemphyxmini0937-satisfiesslidewireforcelaws}
  Eliminating the induced emf and current gives the general drag-force magnitude for this perpendicular slidewire geometry
\end{lemma}
\begin{proof}
  The conclusion follows from the governing relations and figure readouts named in the statement of magnetic force magnitude formula.
\end{proof}
% --- Archon declaration topology end ---
% --- Archon physics formalization source end ---
